% Generated draft: requires mathematical and source review.
\documentclass[11pt]{article}
\usepackage[margin=25mm]{geometry}
\usepackage{fontspec}
\setmainfont{Latin Modern Roman}
\usepackage{amsmath,amssymb,mathtools,unicode-math}
\setmathfont{Latin Modern Math}
\usepackage{enumitem,needspace,xurl,hyperref}
\setlength{\parindent}{0pt}
\setlength{\parskip}{5pt}
\setlength{\emergencystretch}{3em}
\newcommand{\N}{\mathbb{N}}
\newcommand{\Z}{\mathbb{Z}}
\newcommand{\Q}{\mathbb{Q}}
\newcommand{\R}{\mathbb{R}}
\newcommand{\C}{\mathbb{C}}
\newcommand{\F}{\mathbb{F}}
\newcommand{\A}{\mathbb{A}}
\newcommand{\PP}{\mathbb{P}}
\newcommand{\E}{\mathbb{E}}
\newcommand{\eps}{\varepsilon}
\newcommand{\dd}{\,\mathrm{d}}
\begin{document}
\section*{UnsolvedMath ID: 801173}
\textbf{Source problem code:} catalogue entry
\subsection{Definitions and mathematical statement}
Let \(n\geq 1\). In the standard terminology, the Schwartz class is
\[
\mathcal S(\R^n)=\left\{f\in C^\infty(\R^n;\C):\sup_{x\in\R^n}|x^\alpha\partial^\beta f(x)|<\infty\text{ for all }\alpha,\beta\in\Z_{\geq0}^n\right\}.
\]
Here \(x^\alpha=\prod_{i=1}^n x_i^{\alpha_i}\) and \(\partial^\beta=\prod_{i=1}^n\partial_{x_i}^{\beta_i}\). An integer expansive dilation is a matrix \(A\in M_n(\Z)\) whose eigenvalues all have modulus greater than \(1\). The prescribed translation lattice is \(\Gamma=\Z^n\).

The supplied question asks whether Schwartz class wavelets exist for these dilations and this lattice. The record does not define ``wavelet'' or specify the number of generators. Under the finite orthonormal wavelet-family interpretation, a precise candidate formulation is: for every \(n\geq1\) and every integer expansive matrix \(A\in M_n(\Z)\), do there exist a positive integer \(r\) and functions \(\psi_1,\ldots,\psi_r\in\mathcal S(\R^n)\) such that
\[
\left\{x\longmapsto |\det A|^{j/2}\psi_\ell(A^j x-k):j\in\Z,\ k\in\Z^n,\ 1\leq\ell\leq r\right\}
\]
is an orthonormal basis of \(L^2(\R^n;\C)\), with respect to Lebesgue measure?

This is a conditional formulation requiring confirmation against the source's wavelet conventions. The supplied record does not establish whether a single generator, a prescribed number of generators, or another wavelet notion is intended. Its description of the Schwartz class as all ``\(\infty\)-regular'' functions omits the rapid-decay conditions in the standard definition above; smoothness alone is insufficient. No additional restriction on \(A\) or \(n\) is supplied.
\subsection{Short English statement}
The question asks whether every integer expansive dilation, with translations by \(\Z^n\), admits wavelets in the Schwartz class. The displayed formulation uses a finite orthonormal family; the intended wavelet convention and number of generators require source confirmation. The supplied registry labels the problem open.
\subsection{Sources}
\begin{itemize}
\item Supplied local UnsolvedMath registry record: permanent ID \(801173\), exact source problem code \texttt{catalogue entry}, title \emph{Open Problems in Wavelet Theory --- Problem 12}. The registry attributes the problem to Bownik and Rzeszotnik, with proposed year 2019, and supplies the status label ``open''; that status is not independently verified here.
\item Bownik and Rzeszotnik, \emph{Open Problems in Wavelet Theory} (2019), Problem 12. Supplied source URL: \url{https://arxiv.org/abs/1905.08230}. The supplied provenance records access on 2026-10-16 and extraction by \texttt{arxiv\_latex}; these are supplied provenance claims, not a claim of access or verification in preparing this draft.
\end{itemize}
\end{document}
